% A limitation of a proposed route, not a geometric counterexample to A.4.
\subsection{Rational algebraic monodromy does not determine the residual image}

\begin{proposition}[A compact open residual countermodel]
\label{prop:root-rational-residual-obstruction}
\status{Proved by a direct argument; independent audit requested.}
Let \(p\) be a prime, \(d\ge1\), and
\(P\subseteq\operatorname{GL}_d(\mathbf F_p)\) any subgroup.
Write \(L=\mathbf Z_p^d\), and let
\[
 \operatorname{red}:\operatorname{GL}_{\mathbf Z_p}(L)
       \longrightarrow\operatorname{GL}_{\mathbf F_p}(L/pL)
\]
be reduction, a continuous surjective group homomorphism. Then
\(\Gamma=\operatorname{red}^{-1}(P)\), with its natural profinite
topology, acts continuously on \(L\), has residual image exactly
\(P\), and has Zariski closure
\(\operatorname{GL}_{d,\mathbf Q_p}\) on
\(L\otimes_{\mathbf Z_p}\mathbf Q_p\).
In particular its rational representation is absolutely irreducible,
irrespective of reducibility of the residual representation.
\end{proposition}
\begin{proof}
An invertible matrix over \(\mathbf F_p\) admits an entrywise lift
whose determinant is a \(p\)-adic unit; this proves surjectivity of
\(\operatorname{red}\), hence the exact residual image. The inverse
image of a subgroup of a finite discrete group is a compact open
subgroup. It contains \(1+p\operatorname{Mat}_d(\mathbf Z_p)\).

Every polynomial over \(\mathbf Q_p\) vanishing on the latter
set is zero. Indeed substitute \(x_{ij}=\delta_{ij}+p y_{ij}\).
The resulting polynomial vanishes on the Cartesian product
\(\mathbf Z_p^{d^2}\). In one variable infinitely many roots force
it to vanish identically; induction on the variables proves the same
assertion for the product. The substitution is an invertible affine
change over \(\mathbf Q_p\). A regular function on
\(\operatorname{GL}_{d,\mathbf Q_p}\) is a polynomial divided by
a power of the determinant, so the same reasoning proves the stated
Zariski density. The argument remains true after extending scalars to
an algebraic closure: expand the finitely many coefficients in a
\(\mathbf Q_p\)-basis of their span and apply the preceding argument
to each coefficient polynomial. Any invariant subspace in the
algebraic closure is consequently invariant under the full general
linear group, whose natural representation is irreducible.
\end{proof}

\begin{example}[The absence of fixed vectors and covectors still permits a parabolic]
\label{ex:root-rational-residual-parabolic}
\status{Proved by a direct argument and an exact finite supplement; independent audit requested.}
For \(r>s\ge2\), put \(d=r+s\) and
\[
 P=\left\{
  \begin{pmatrix} A&B\\0&D\end{pmatrix}:
  A\in\operatorname{GL}_r(\mathbf F_p),\quad
  D\in\operatorname{GL}_s(\mathbf F_p),\quad
  B\in\operatorname{Mat}_{r\times s}(\mathbf F_p)
 \right\}.
\]
This is a proper subgroup preserving the proper nonzero subspace
\(U=\mathbf F_p^r\oplus0\). For the natural module
\(V=\mathbf F_p^d\), nevertheless
\(V^P=0=(V^\vee)^P\), and \(P\) contains the full natural
\(\operatorname{GL}_r(\mathbf F_p)\) block with
\(r>d/2\). Its inverse image in
\(\operatorname{GL}_d(\mathbf Z_p)\) has the full rational
algebraic monodromy of the preceding proposition.
\end{example}
\begin{proof}
Block multiplication and inversion preserve the displayed form, so
it is a subgroup. An elementary matrix carrying one vector of \(U\)
to a vector with nonzero second component does not preserve \(U\),
proving properness. The subgroup of block diagonal matrices acts
independently on the two summands. For any \(m\ge2\), the natural
\(\operatorname{GL}_m(\mathbf F_p)\)-module has no fixed nonzero
vector: a nonzero coordinate can be added to a different coordinate
by an elementary transvection. The same argument after transpose
applies to covectors. Apply this separately to both diagonal blocks.
The large natural block is the subgroup with \(D=1,B=0\), and
the final assertion follows from
Proposition~\ref{prop:root-rational-residual-obstruction}.

The retained program
\path{code/check_rational_residual_parabolic.py}
checks the explicit finite instance \(p=2,r=3,s=2\). It enumerates
the full displayed subgroup, checks every member's rank and invariant
subspace, and checks the fixed vectors and covectors against elementary
generators. Its output states exactly those finite claims. The
infinite \(p\)-adic density and all-rank assertions are proved above.
\end{proof}

\paragraph{Scope of the literature route.}
We examined the original preprint of D'Addezio
\cite[Theorem~4.4.12, printed pp.~23--24;
Proposition~5.1.4 and Theorem~5.1.6, printed p.~25]
{DAddezio2023Parabolicity}. The final proofs at these locators were
read, but their deeper dependencies were not all independently
audited. The statements concern rational isocrystals and rational
algebraic monodromy; Theorem~5.1.6 in particular states rational
semisimplicity for an abelian scheme of constant Newton slopes.
We do not import these results as a residual-image theorem. Even
full rational algebraic monodromy and absence of residual fixed
vectors and covectors leave the explicit parabolic above possible.
Additional integral or residual geometry is necessary for A.4.
This example is an abstract continuous representation; it does not
assert its realization by the hypersurface family in the question,
and it does not shrink the remaining geometric main gate.
